\begin{abstract}
Large-scale unsupervised language models (LMs) are capable of acquiring extensive world knowledge and exhibit certain reasoning abilities, yet directing their outputs with accuracy is challenging because their training lacks explicit supervision. To address this, current approaches acquire human feedback indicating preferences between pairs of model outputs, and then fine-tune these unsupervised LMs so that their behavior aligns more closely with the collected preferences, typically employing reinforcement learning from human feedback (RLHF). Despite its effectiveness, RLHF involves a multi-stage and potentially unstable pipeline, in which a reward model is first trained to predict human preferences, and then the unsupervised LM is further adapted using reinforcement learning, maximizing this reward without diverging excessively from the initial model parameters. In this work, we propose a novel reparameterization of the reward model within RLHF, making it possible to directly recover the optimal policy in closed form. This insight reduces the standard RLHF objective to a straightforward classification loss, leading to an algorithm we term \textit{Direct Preference Optimization} (DPO). DPO is characterized by stability, strong empirical performance, and low computational demands, and does not require sampling from the language model or extensive hyperparameter search during the fine-tuning process. Our results demonstrate that DPO aligns LMs to human preferences as effectively or better than established techniques. In particular, DPO outperforms PPO-based RLHF in controlling the sentiment of generated text, and matches or exceeds the quality of model outputs in summarization and single-turn conversation tasks, while offering a more streamlined and accessible training method.
\end{abstract}

\section{Introduction}
Large-scale unsupervised language models (LMs), when trained on extensive datasets, often develop unexpected proficiencies~\citep{chowdhery2022palm, brown2020language, touvron2023llama, bubeck2023sparks}. Nevertheless, such models are trained on corpora created by humans who differ widely in their aims, standards, and expertise. Inevitably, some objectives and competencies represented in this data are not desirable to replicate; for instance, we might prefer our AI-based code assistant to \textit{recognize} prevalent coding errors so it can correct them, yet we wish it to emulate the exceptional, albeit infrequent, high-level programming skills embedded in the training corpus when actually producing code. Likewise, while it is beneficial for a language model to be \textit{cognizant} of popular misconceptions that might be accepted by half the population, it is critical that the model does not reproduce these misconceptions in half of its responses to relevant queries. Hence, identifying and amplifying the model's \emph{preferred outputs and behaviors} from among its diverse \textit{knowledge and skillset} is essential for constructing AI systems that are robust, effective, and amenable to control~\citep{ouyang2022training}. Standard approaches currently guide LMs to align with human values via reinforcement learning (RL) based preference modeling; this paper demonstrates, however, that the RL-driven objective utilized by these techniques is precisely equivalent to optimizing a simple binary cross-entropy loss, thus streamlining the preference learning process significantly.

Broadly, present methodologies embed desired model behaviors through curated datasets reflecting human preferences regarding safe and beneficial conduct. This preference refinement process is typically carried out following an initial phase of unsupervised pre-training on large textual datasets. The most direct method for incorporating preferences is supervised fine-tuning using annotated examples of excellent responses, but the predominant approach in practice is reinforcement learning from human or AI feedback (RLHF/RLAIF; \citep{christiano2017deep,bai2022constitutional}). RLHF operates by constructing a reward model from human-annotated preference data and subsequently applying RL techniques to tune the language model, maximizing the likelihood of producing outputs that score highly on this reward model, while preventing substantial deviation from the initial policy. Although RLHF is notable for achieving advanced conversational and coding performances, its overall pipeline is much more intricate than straightforward supervised learning, since it involves the development of several distinct LMs and requires frequent policy sampling during training, all of which contribute to considerable computational resource demands.

In this work, we demonstrate a method for aligning a language model with human preferences without relying on explicit reward models or reinforcement learning techniques. We introduce \textit{{Direct Preference Optimization} (DPO)}, a method that, while implicitly optimizing the same target as traditional RLHF approaches (maximizing reward subject to a KL-divergence constraint), offers greater simplicity in both implementation and training procedures. The central idea behind DPO is to adjust the model such that it enhances the log-odds of preferred responses over less desirable alternatives. Crucially, DPO incorporates a dynamic importance weighting on a per-sample basis, mitigating model collapse issues observed when utilizing a straightforward probability ratio objective. DPO, akin to previous algorithms, employs a formal preference model (e.g., the Bradley-Terry model; \cite{bradley1952rankanalysis}) to assess the alignment between a proposed reward function and actual preference judgments. Nevertheless, while prior techniques apply this preference model to generate a loss for fitting a reward model, followed by policy optimization on the learned reward, DPO reformulates the preference loss directly in terms of the policy through a variable substitution. Consequently, given annotated human preference data on model outputs, DPO is able to directly optimize a policy by minimizing a binary cross-entropy loss, effectively yielding the optimal policy for an implicit reward function shaped by preference feedback.

The principal innovation of this paper is Direct Preference Optimization (DPO), a streamlined, RL-free approach for optimizing language models according to preference data. Empirical results indicate that DPO matches or exceeds the effectiveness of prevailing techniques, such as PPO-based RLHF, across a variety of preference-driven applications—including sentiment transformation, summarization, and conversational modeling—in models scaling up to 6B parameters.

\section{Related Work}

Large-scale self-supervised language models have demonstrated the capacity to solve certain problems without prior examples (zero-shot; \citep{radford2019language}) or with minimal prompt examples (few-shot; \citep{gpt3,megatron,chowdhery2022palm}). Nonetheless, their efficacy on specific applications and their ability to follow user instructions can be substantially heightened through fine-tuning on corpora composed of explicit instructions and responses crafted by humans \citep{mishra-etal-2022-cross,sanh2022multitask,chung2022scaling,thoppilan2022lamda}. This process, commonly referred to as ‘instruction-tuning’, enhances LLMs’ adaptability to novel instructions beyond their fine-tuning set and generally makes them more practical for users \citep{chung2022scaling}. Although instruction-tuning has shown strong empirical success, collecting relative assessments of output quality from non-experts is typically more scalable than assembling expert-generated demonstrations. Consequently, subsequent research has refined LLMs using datasets composed of human preference judgments, leading to advances in tasks such as translation \citep{kreutzer-etal-2018-reliability}, summarization \citep{stiennon2022learning,ziegler2020finetuning}, narrative generation \citep{ziegler2020finetuning}, and following instructions \citep{ouyang2022training,ramamurthy2023is}. The prevalent paradigm first entails training a reward model—typically parameterized by a neural network—to predict human preferences, often leveraging frameworks like the Bradley-Terry model \citep{bradley1952rankanalysis}. Subsequently, the language model is fine-tuned to maximize these learned rewards using reinforcement learning strategies such as REINFORCE \citep{williams1992reinforce}, proximal policy optimization (PPO; \cite{schulman2017proximal}), or related techniques \citep{ramamurthy2023is}. In parallel, another strand of research employs LLMs trained for instruction adherence with human feedback to synthesize additional preference data targeting qualities like safety or non-harmfulness \citep{bai2022constitutional}, relying on weak supervision signals—such as rubric-based textual guidelines—for annotating the generated data. This methodological convergence blends developments from two principal areas: reinforcement learning applications to language model training for diverse objectives~\citep{Ranzato2015SequenceLT,paulus2018a,wu2018learning}, and broader strategies for leveraging human preference information \citep{christiano2017deep,kupcsik2018learning}. However, the implementation of reinforcement learning for adapting large language models based on relative human feedback continues to pose considerable difficulties in practice; this study introduces a theoretically grounded method for optimizing over relative preferences that does not rely on RL.

Beyond linguistic applications, the problem of policy learning from preferences has been explored within both the bandit and reinforcement learning paradigms, leading to a variety of methodological proposals. The contextual dueling bandit (CDB) problem, for instance, generalizes contextual bandit learning by replacing explicit reward signals with preferences or rankings over actions \citep{yue2012karmed,dudik2015contextual}. When absolute rewards are not available, analysis of CDBs introduces the concept of a \textit{von Neumann winner}: a policy whose expected rate of outperforming any competing policy is no less than 50\% \citep{dudik2015contextual}. It is noteworthy, however, that CDBs are characterized by the sequential acquisition of preference feedback, whereas in typical settings involving human preference learning, we often operate on a static, pre-collected batch of annotated action pairs \citep{yan2022human}. In a related domain, \textit{preference-based RL} (PbRL) focuses on utilizing binary preference feedback—originating from an unknown 'scoring' function—as a surrogate for direct reward feedback \citep{BusaFekete2014,ruiz2023dueling}. The spectrum of PbRL algorithms encompasses approaches that support off-policy reuse of preferences, yet these are usually structured to first infer the underlying scoring or reward function, followed by policy optimization based on this estimate \citep{jain2013learning,BusaFekete2014,christiano2017deep,sadigh2017active,kupcsik2018learning}. Contrary to this two-stage paradigm, we propose an integrated method that optimizes the policy directly to conform with given preference data.

\section{Preliminaries}\label{section:prelims}

Here, we provide an overview of the RLHF framework as described by \citeauthor{ziegler2020finetuning}, as well as its subsequent iterations \citep{stiennon2022learning, bai2022training, ouyang2022training}. This framework typically encompasses three core steps: (1) supervised fine-tuning (SFT), (2) collection of preference data and subsequent reward model estimation, and (3) reinforcement learning-based policy refinement.

\textbf{SFT}: The RLHF workflow begins by leveraging supervised learning to further train a large language model (LM) on a curated set of high-quality examples relevant to target applications (e.g., conversation, summarization), yielding the fine-tuned model $\pi^\text{SFT}$.

\textbf{Reward Modelling Phase}: The second stage involves using the SFT model to generate response pairs to prompts $x$, resulting in samples $(y_1, y_2)\sim \pi^\text{SFT}(y \mid x)$. Human annotators then review these pairs and indicate their preference, expressed as $y_w\succ y_l \mid x$, where $y_w$ and $y_l$ identify the more and less favored completions, respectively. Although the preferences originate from an unobservable underlying reward function $r^*(y, x)$, we approximate this relationship using probabilistic models. One of the most widely adopted models for representing preferences is the Bradley-Terry (BT) framework \cite{bradley1952rankanalysis} (although the more general Plackett-Luce model \citep{plackett1975analysis, luce2012individual} is suitable when full rankings of outputs are available). Under the BT model, the true human preference distribution $p^*$ can be formulated as:

\begin{equation}\label{eq:bradley-terry}
    p^*(y_1\succ y_2 \mid x)=\frac{\exp\left(r^*(x, y_1)\right)}{\exp\left(r^*(x, y_1)\right) + \exp\left(r^*(x, y_2)\right)}.
\end{equation}

Given a fixed dataset of pairwise preference comparisons $\mathcal{D}=\bigl\{x^{(i)}, y_w^{(i)}, y_l^{(i)}\bigr\}_{i=1}^N$ drawn from $p^*$, we can represent the reward function by a parameterized model $r_{\phi}(x, y)$, and optimize its parameters by maximizing the likelihood. Interpreting the learning process as binary classification, the associated loss is the negative log-likelihood:

\begin{equation}\label{eq:reward_model}
    \mathcal{L}_R(r_{\phi}, \mathcal{D}) = -\mathbb{E}_{(x, y_w, y_l)\sim \mathcal{D}}\bigl[\log \sigma(r_{\phi}(x, y_w)- r_{\phi}(x, y_l))\bigr]
\end{equation}

where $\sigma$ denotes the logistic sigmoid. For large language models, $r_{\phi}(x, y)$ is commonly initialized from the SFT model $\pi^\text{SFT}(y \mid x)$, augmented with a linear projection atop the final transformer layer to yield a scalar reward prediction \cite{ziegler2020finetuning}. To reduce the variance of the learned reward, it is standard practice to normalize so that $\mathbb{E}_{x,y\sim \mathcal{D}}\left[r_\phi(x, y)\right] = 0$ for any $x$.

\textbf{RL Fine-Tuning Phase}: In the subsequent RL stage, the acquired reward model serves to guide the training of the language model. Following established methods~\citep{jaques2017sequence, jaques2020human}, policy optimization is posed as

\begin{equation}\label{eq:RL}
\max_{\pi_{\theta}}  \mathbb{E}_{x\sim \mathcal{D}, y\sim \pi_{\theta}(y \mid x)}\bigl[r_{\phi}(x, y)\bigr] - \beta\mathbb{D}_{\textrm{KL}}\bigl[\pi_{\theta}(y\mid x)\mid \mid \pi_\text{ref}(y\mid x)\bigr],
\end{equation}

where the hyperparameter $\beta$ moderates divergence from a reference policy $\pi_\text{ref}$—specifically, the initial SFT model $\pi^\text{SFT}$. The policy $\pi_\theta$ is likewise initialized from $\pi^\text{SFT}$ in practice. Imposing this constraint curbs excessive deviation from the distribution under which the reward model remains reliable, thus helping retain answer diversity and mitigating collapse to a few high-scoring responses. As the process of generating language is inherently discrete, direct differentiation through this objective is infeasible, necessitating reinforcement learning optimization methods. A common practice \citep{ziegler2020finetuning, stiennon2022learning, bai2022training, ouyang2022training} is to define the reward as ${r(x, y) = r_{\phi}(x, y) -\beta (\log \pi_{\theta}(y\mid x) - \log \pi_\text{ref}(y\mid x))}$ and employ PPO \cite{schulman2017proximal} to maximize it.

\section{Direct Preference Optimization}\label{sec:DPO}

Recognizing the limitations of applying reinforcement learning techniques for large-scale language model fine-tuning, our objective is to develop a streamlined method for optimizing policies directly from preferences. Departing from conventional RLHF pipelines, which require learning a reward model before RL-based optimization, we adopt a specialized reward parameterization that allows the optimal policy to be determined in closed form, bypassing the need for RL altogether. As elaborated in the following sections, our main contribution is to exploit a direct analytic relationship between reward functions and their corresponding optimal policies. This insight enables reparameterization: the loss on reward functions can be recast as a loss over policies themselves. This change of variables circumvents training an explicit reward network while retaining alignment with standard models of human preference, such as the Bradley-Terry formulation. Effectively, the policy network doubles as both the generator and an implicit reward estimator.

\textbf{Derivation of the DPO objective.} Consider the standard RL objective, as stated in Eq.~\ref{eq:RL}, with an arbitrary reward function $r$. Building on previous studies~\citep{peters2007reinforcement, peng2019advantage, korbak2022reinforcement, go2023aligning}, it follows that the optimal policy for the KL-regularized reward maximization problem in Eq.~\ref{eq:RL} has the structure:

\begin{equation}\label{eq:op_policy}
    \pi_r(y\mid x) = \frac{1}{Z(x)}\pi_\text{ref}(y\mid x)\exp\left(\frac{1}{\beta}r(x, y)\right),
\end{equation}

with $Z(x) =\sum_{y}\pi_\text{ref}(y\mid x)\exp\left(\frac{1}{\beta}r(x, y)\right)$ representing the normalization constant. Detailed derivation can be found in Appendix \ref{app:derivation1}. In practical applications, even when replacing $r$ by its MLE estimate $r_{\phi}$ for the true reward $r^*$, evaluating the partition function $Z(x)$ remains computationally intensive~\citep{korbak2022reinforcement, go2023aligning}, making direct use of this formulation challenging. Nevertheless, Eq.~\ref{eq:op_policy} can be rearranged to isolate the reward function as a function of the optimal policy $\pi_r$, the reference policy $\pi_\text{ref}$, and the (possibly intractable) partition function $Z(\cdot)$. In particular, applying the logarithm to both sides of Eq.~\ref{eq:op_policy} and rearranging yields:

\begin{equation}\label{eq:main_eq}
    r(x,y) =\beta \log \frac{\pi_r(y\mid x)}{\pi_\text{ref}(y\mid x)} + \beta \log Z(x).
\end{equation}

This transformation applies equally to the ground-truth reward $r^*$ and its associated optimal policy $\pi^*$. Importantly, the Bradley-Terry model concerns itself with reward differences between two outcomes, namely, ${p^*(y_1 \succ y_2 \mid x) = \sigma(r^*(x, y_1) - r^*(x, y_2))}$. Substituting the reparameterized form from Eq.~\ref{eq:main_eq} for $r^*(x, y)$ into the preference expression in Eq.~\ref{eq:bradley-terry} eliminates the partition function, thus expressing human preference probability solely as a function of the optimal policy $\pi^*$ and the reference policy $\pi_\text{ref}$. Therefore, under the Bradley-Terry framework, the optimal RLHF policy $\pi^*$ abides by the following preference equation:

\begin{equation}\label{eq:objective}
    p^*(y_1\succ y_2 \mid x)=\frac{1}{1 + \exp\left(\beta \log \frac{\pi^*(y_2\mid x)}{\pi_\text{ref}(y_2\mid x)} - \beta \log \frac{\pi^*(y_1\mid x)}{\pi_\text{ref}(y_1\mid x)}\right)}
\end{equation}

Further details are available in Appendix~\ref{app:derivation2}. While Eq.~\ref{eq:objective} is specifically derived with the Bradley-Terry preference model, analogous derivations using the broader class of Plackett-Luce models~\citep{plackett1975analysis, luce2012individual} are provided in Appendix~\ref{app:plackett_luce_models}.

Having related the likelihood of human preference data to the optimal policy rather than the explicit reward, it becomes possible to specify a maximum likelihood criterion for a parameterized policy $\pi_\theta$. Following the reward modeling perspective (cf. Eq.~\ref{eq:reward_model}), the corresponding policy objective is:

\begin{equation}\label{eq:optimum_model}
    \mathcal{L}_\text{DPO}(\pi_{\theta}; \pi_\text{ref}) = -\mathbb{E}_{(x, y_w, y_l)\sim \mathcal{D}}\left[\log \sigma \left(\beta \log \frac{\pi_{\theta}(y_w\mid x)}{\pi_\text{ref}(y_w\mid x)} - \beta \log \frac{\

In this formulation, we learn an implicit reward function through an alternative parameterization, resulting in an optimal policy that coincides with $\pi_\theta$. Additionally, as this method aligns with a reparametrized Bradley-Terry model, it inherits certain theoretical guarantees, including consistency under appropriate conditions on the distribution of preference data \cite{bong2022generalized}. We elaborate on further theoretical aspects of DPO and its relation to other approaches in Section~\ref{sec:theory}.

\textbf{Interpretation of the DPO update.} To gain insight into the workings of DPO, it is instructive to examine the gradient of the objective function $\mathcal{L}_\text{DPO}$. Specifically, the gradient with respect to the model parameters $\theta$ takes the following form:

\begin{multline*}\label{eq:gradient}
    \nabla_\theta \mathcal{L}_\text{DPO}(\pi_\theta;\pi_\text{ref}) = \\ -\beta\mathbb{E}_{(x, y_w, y_l) \sim \mathcal{D}} \bigg[\underbrace{\sigma(\hat{r}_\theta(x, y_l) - \hat{r}_\theta (x, y_w))}_\text{higher weight when reward estimate is wrong}\bigg[\underbrace{\nabla_\theta\log \pi(y_w \mid x)}_\text{increase likelihood of $y_w$} - \underbrace{\nabla_\theta\log\pi(y_l \mid x)}_\text{decrease likelihood of $y_l$}\bigg]\bigg],
\end{multline*}

where the implicit reward, defined as $\hat{r}_\theta(x, y) = \beta \log \frac{\pi_\theta(y \mid x)}{\pi_\text{ref}(y \mid x)}$, is parameterized by both the current model $\pi_\theta$ and the reference model $\pi_\text{ref}$ (for additional details, see Section~\ref{sec:theory}). Intuitively, this gradient acts to make the preferred outputs $y_w$ more probable and the less preferred ones $y_l$ less probable under $\pi_\theta$. Crucially, each training instance is weighted by the extent to which the implicit reward $\hat{r}_\theta$ mistakenly assigns a higher score to the less preferred outcome $y_l$—the scale of this error is determined by $\beta$, which also reflects the strength of the KL regularization. Empirically, this form of importance weighting is significant: a simpler variant that omits the weighting factor can lead to undesirable language model behavior, as shown in Appendix Table~\ref{tab:unlikelihood_generations}.

\textbf{DPO overview.} The typical DPO workflow comprises the following steps: 1) For each prompt $x$, generate completions $y_1, y_2$ from $\pi_\text{ref}(\cdot \mid x)$, and then annotate these pairs based on human preference to construct the offline preference dataset $\mathcal{D} = \{x^{(i)}, y_w^{(i)}, y_l^{(i)}\}_{i=1}^N$; 2) train the language model $\pi_\theta$ by minimizing $\mathcal{L}_\text{DPO}$ using the specified $\pi_\text{ref}$, dataset $\mathcal{D}$, and a chosen $\beta$. 
In actual applications, rather than creating new samples and collecting human preference labels, practitioners often make use of available public preference datasets. Since these datasets are typically sampled using $\pi^\text{SFT}$, we set $\pi_\text{ref} = \pi^\text{SFT}$ when it is accessible. If $\pi^\text{SFT}$ is not accessible, we instead initialize $\pi_\text{ref}$ by maximizing the likelihood over the preferred completions ${(x, y_w)}$, formally, ${\pi_\text{ref} = \argmax_{\pi}\mathbb{E}_{x, y_w \sim \mathcal{D}}\left[\log \pi(y_w \mid x)\right]}$. This approach helps to reduce the distribution gap between the unavailable true reference distribution and the $\pi_\text{ref}$ utilized in DPO. Additional implementation specifics and hyperparameter selections are detailed in Appendix~\ref{app:implementation}.

\section{Theoretical Analysis of DPO} This section presents additional insights into the DPO framework, offering theoretical justification and discussing its strengths relative to actor-critic methods commonly used in RLHF (like PPO~\cite{schulman2017proximal}).

\label{sec:theory}

\subsection{Your Language Model Is Secretly a Reward Model} DPO achieves policy learning by unifying the process of reward modeling and RL, allowing both to be handled through a single maximum likelihood formulation. The optimization criterion in Eq. \ref{eq:main_eq} can be seen as equivalent to the Bradley-Terry model, where the reward is parameterized as $r^*(x, y) = \beta \log\frac{\pi^*_\theta(y \mid x)}{\pi_\text{ref}(y \mid x)}$, and optimizing the parametric model $\pi_{\theta}$ matches the procedure for reward model fitting in Eq. \ref{eq:reward_model} after a suitable reparameterization. In this section, we develop the theoretical foundations for this reparameterization, demonstrate that it does not limit the expressiveness of the reward model class, and establish that it is sufficient for exactly recovering the optimal policy. Our analysis begins by introducing an equivalence relation over reward functions.

\begin{definition}
We define two reward functions $r(x, y)$ and $r'(x, y)$ as equivalent if and only if there exists a function $f$ such that ${r(x, y)-r'(x, y) = f(x)}$.     
\end{definition}

It is straightforward to verify that this defines an equivalence relation, which divides the set of reward functions into disjoint equivalence classes. This leads to the following two lemmas:

\begin{lemma}\label{lemma:same_prefrence}
Within the Plackett-Luce preference model—and by extension, the Bradley-Terry model—any two reward functions belonging to the same equivalence class result in identical preference distributions.
\end{lemma}

\begin{lemma}\label{lemma:same_policy}
    Any pair of reward functions from a given equivalence class will generate the same optimal policy when solving the associated constrained RL problem.
\end{lemma}

The arguments are immediate and can be found in Appendix \ref{app:lemma1}. The first lemma points to a known under-specification characteristic inherent to the Plackett-Luce model family \cite{plackett1975analysis}. This under-determination typically necessitates supplementary identifiability conditions to secure guarantees regarding the MLE solutions of Eq. \ref{eq:reward_model} \cite{bong2022generalized}. The second lemma establishes that, with respect to finding the optimal policy, the specific representative chosen within an equivalence class of reward functions is inconsequential; thus, it suffices to recover any element of the optimal equivalence class. The following Theorem, proved in Appendix~\ref{app:thm1}, formalizes this notion:

\begin{theorem}\label{thm:main}
    Assuming mild technical conditions, every equivalence class of rewards compatible with the Plackett-Luce (and in particular, Bradley-Terry) frameworks admits a representation via the reparameterization ${r(x, y) = \beta \log \frac{\pi(y\mid x)}{\pi_\text{ref}(y\mid x)}}$, where $\pi(y\mid x)$ is a model and $\pi_\text{ref}(y \mid x)$ is a specified reference model.
\end{theorem}

\begin{sproof}
    Let us take any reward function $r(x, y)$, and denote the associated optimal policy by $\pi_r(y \mid x)$, defined as in Eq. \ref{eq:op_policy}. We wish to show that a reward function within the equivalence class of $r$ can be written in the form stated above. Define the projection operator $f$ by  
\begin{equation}
    f(r; \pi_\text{ref}, \beta)(x, y) = r(x, y) - \beta\log\sum_{y}\pi_\text{ref}(y\mid x)\exp\left(\frac{1}{\beta}r(x, y)\right)
\end{equation}
Here, $f$ subtracts the log of the partition function (computed using $\pi_\text{ref}$) from $r(x, y)$. Since this normalization term only depends on $x$, $f(r; \pi_\text{ref}, \beta)(x, y)$ differs from $r(x, y)$ by a function of $x$ alone, thus ensuring it remains within the equivalence class of $r(x, y)$. Substituting the right-hand side of Eq.~\ref{eq:main_eq} for $r$, we observe that $f(r; \pi_\text{ref}, \beta)(x, y) = \beta \log \frac{\pi_r(y\mid x)}{\pi_\text{ref}(y\mid x)}$. Therefore, this projection produces an equivalence-class representative in the required parameterization, confirming that this reparameterization encompasses all possible reward functions in the class.
\end{sproof}

Theorem~\ref{thm:main} can also be interpreted as identifying the particular reward function chosen by the DPO reparameterization within each equivalence class. Specifically, it selects the reward function that fulfills:

\begin{equation}\label{eq:lag_p}
     \sum_{y}\underbrace{\pi_\text{ref}(y\mid x)\exp\left(\frac{1}{\beta}r(x, y)\right)}_{=\pi(y\mid x)\text{, using Thm.~\ref{thm:main} reparam.}} = 1,
\end{equation}

meaning that $\pi(y\mid x)$ constitutes a legitimate probability distribution (i.e., its values are non-negative and sum to one). By referencing Eq.~\ref{eq:op_policy}, one observes that Eq.~\ref{eq:lag_p} represents the partition function corresponding to the optimal policy determined by the reward function $r(x, y)$. The central contribution of the DPO algorithm is its capacity to enforce specific constraints within the otherwise under-determined Plackett-Luce family of preference models (notably including the Bradley-Terry model), thereby maintaining the expressive power of the reward models, while ensuring that the optimal policy in Eq.~\ref{eq:op_policy} admits a closed-form expression for all $x$.

\subsection{Instability of Actor-Critic Algorithms} Our framework also facilitates the identification of instability issues prevalent in actor-critic algorithms like PPO, frequently employed during the RLHF process. Concentrating on the RL fine-tuning procedure discussed in Section~\ref{section:prelims}, we can relate this scenario to the control as inference perspective \cite{levine2018reinforcement} for the constrained RL objective posed in Eq.~\ref{eq:RL}. By considering a parameterized policy $\pi_{\theta}(y\mid x)$, the objective is to minimize $\mathbb{D}_{\text{KL}}[\pi_{\theta}(y|x) \mid \mid \pi^*(y\mid x)]$, where $\pi^*$, derived in Eq.~\ref{eq:optimum_model}, represents the optimal policy corresponding to reward $r_{\phi}(y, x)$. Some manipulations yield the following maximization objective:

\begin{equation}\label{eq:AC}
    \max_{\pi_{\theta}}\mathbb{E}_{\pi_{\theta}(y\mid x)}\bigg[\underbrace{r_{\phi}(x, y) -\beta\log\sum_{y}\pi_\text{ref}(y\mid x)\exp\left(\frac{1}{\beta}r_{\phi}(x, y)\right)}_{f(r_{\phi}, \pi_\text{ref}, \beta)} - \underbrace{\beta\log\frac{\pi_{\theta}(y\mid x)}{\pi_\text{ref}(y\mid x)}}_{\text{KL}}\bigg]
\end{equation}

The objective under consideration aligns with those targeted in earlier research 
\citep{ziegler2020finetuning, stiennon2022learning, bai2022training, ouyang2022training}, employing a DPO-compatible reward within the reward class $r_{\phi}$. Within this framework, the normalization factor present in $f(r_{\phi}, \pi_\text{ref}, \beta)$ can be understood as representing the soft value function corresponding to the reference policy $\pi_\text{ref}$. Although this normalization does not alter the location of the optimum, omitting it may result in a policy gradient estimator with increased variance, thereby impeding stable learning. Employing a learned value function to address this normalization is possible, but it introduces optimization difficulties of its own. Previous studies have circumvented this by utilizing a baseline derived from human completions—essentially estimating the normalizing term via a single Monte-Carlo sample. In contrast, the DPO reparameterization generates a reward structure that obviates the need for any such baselines.

\section{Experiments}
Here, we undertake an empirical analysis of DPO's efficacy in training policies from preference data. Our initial investigation, conducted in a controlled environment focused on text generation, aims to assess DPO's proficiency in balancing reward maximization with the constraint of KL-divergence relative to a reference policy, especially in comparison to widely-used preference optimization methods like PPO. Subsequently, we scale our evaluation to larger model architectures and tackle more challenging RLHF scenarios, encompassing tasks like summarization and dialog generation. The empirical results suggest that, with minimal hyperparameter tuning, DPO generally matches or surpasses established benchmarks such as RLHF with PPO and surpasses strategies involving selection among $N$ generated samples scored by a learned reward function. The subsequent sections detail the experimental methodology, with further specifics provided in Appendix~\ref{app:exp_details}.

\textbf{Tasks.} We investigate three distinct open-ended text generation tasks in our experimental setup. In all settings, algorithms are tasked with learning a policy using a preference dataset $\mathcal{D}=\bigl\{x^{(i)}, y_w^{(i)}, y_l^{(i)}\bigr\}_{i=1}^N$. The first task, \textbf{controlled sentiment generation}, uses $x$ as an initial segment of a movie review sourced from the IMDb dataset \cite{maas-EtAl:2011:ACL-HLT2011}. Here, the model's objective is to generate a continuation $y$ exhibiting positive sentiment. For objective evaluation, we automatically construct preference pairs from generations based on a pre-trained sentiment classifier, ensuring $p(\text{positive}\mid x,y_w)>p(\text{positive}\mid x,y_l)$. For the SFT baseline, GPT-2-large is fine-tuned until convergence using review texts from the IMDB training split (see further details in App~\ref{app:sentiment_details}). The second task, \textbf{summarization}, presents $x$ as a Reddit forum post, with the goal of producing a summary $y$ encapsulating its main ideas. Building on previous studies, we employ the Reddit TL;DR dataset \citep{volske-etal-2017-tl} augmented with human preference annotations from \citeauthor{stiennon2022learning}. The SFT model is derived by further fine-tuning on human-constructed Reddit summaries\footnote{\url{https://huggingface.co/CarperAI/openai_summarize_tldr_sft}} utilizing the TRLX \citep{leandro_von_werra_2023_7790115} system for RLHF. The referenced human preference labels were originally collected by \citeauthor{stiennon2022learning} from outputs generated by a separate but similarly fine-tuned SFT model. The final task, \textbf{single-turn dialogue}, defines $x$ as a user input—ranging from scientific questions to requests for advice—where the policy is expected to generate a response $y$ that is both helpful and engaging. For this, we leverage the Anthropic Helpful and Harmless dialogue dataset \citep{bai2022training}, which comprises 170,000 human-assistant dialogues, with each concluding in two automated responses and a human-annotated preference indicating the superior reply. Since a dedicated SFT model is not provided for this scenario, we develop one by fine-tuning a general-purpose language model exclusively on responses marked as preferred.

\textbf{Evaluation.} Our evaluation strategy employs two distinct methodologies. In the controlled sentiment generation context, where the ground-truth reward function (sentiment classifier) is available, we quantify each algorithm’s performance based on the spectrum it establishes between achieved reward and KL-divergence relative to the reference policy; the reward-KL tradeoff curve, or frontier, can thus be accurately determined. Conversely, for practical applications where the reward function is not directly accessible, we measure algorithmic performance by calculating the \textit{win rate} against a baseline policy, relying on GPT-4 to serve as a surrogate for human judgments. This approach is applied to assess both the quality of summaries and the helpfulness of responses in summarization and single-turn dialogue settings, respectively. For the summarization task, the comparison baseline is the reference summary present in the test set, while in dialogue tasks, the preferred human response from the test set serves as the baseline. Although previous research indicates that LMs often provide more consistent automatic evaluations than conventional metrics \citep{Chen2023ExploringTU}, we conduct a supplementary human subject study in Sec.~\ref{sec:human-judgments} to support our adoption of GPT-4 as an evaluation benchmark. Our results demonstrate that GPT-4’s assessments exhibit high correlation with human judgments, with rates of human-GPT-4 agreement often matching or surpassing those observed among human annotators.

\textbf{Methods.} Alongside DPO, we investigate a variety of prevailing methods for aligning language models with human preferences. For summarization, we use zero-shot prompting with \textbf{GPT-J} \citep{gpt-j}, while for dialogue, we utilize two-shot prompting with \textbf{Pythia-2.8B} \citep{biderman2023pythia}. We additionally consider the \textbf{SFT} approach and introduce the \textbf{Preferred-FT} model, which undergoes supervised finetuning on preferred completions $y_w$ derived from either the SFT model (for sentiment and summarization) or a standard LM (for dialogue). The \textbf{Unlikelihood}~\citep{welleck2019neural} strategy—another form of pseudo-supervised learning—involves maximizing the model’s likelihood on $y_w$ while explicitly \textit{reducing} likelihood on $y_l$; here, the 'unlikelihood' penalty term can be modulated by an optional parameter $\alpha \in [0,1]$. Furthermore, we assess \textbf{PPO} \citep{schulman2017proximal}, which is trained with a reward model learned from preference data, and \textbf{PPO-GT}, an oracle variant that learns directly from the true reward function, applicable in controlled sentiment experiments. For PPO-GT, our sentiment tasks leverage both a standard implementation \cite{leandro_von_werra_2023_7790115} and a modified version featuring normalized rewards and enhanced hyperparameter tuning; these adjustments are also made for standard PPO with learned rewards to boost efficacy. Lastly, we incorporate the \textbf{Best of $N$} approach as a baseline. This method samples $N$ candidate outputs from the SFT (or Preferred-FT in dialogue) model and selects the response that receives the highest score according to a reward model trained on preference annotations. While this baseline is notably effective and sidesteps direct coupling between reward modeling and PPO optimization, it is computationally demanding—even for moderate values of $N$—since every test instance requires generating and evaluating $N$ completions.

\subsection{How well can DPO optimize the RLHF objective?}

In standard RLHF methodologies, the KL-constrained reward maximization objective serves to strike a balance between maximizing reward and keeping the learned policy close to the reference policy. Thus, proper assessment of these algorithms necessitates evaluating both the realized reward and the degree of KL divergence; marginal reward improvements that come at the expense of substantial increases in KL are not inherently advantageous. Figure~\ref{fig:frontier-tldr-main} presents the reward-KL trade-off frontiers for several algorithms within the sentiment analysis framework. For each algorithm, we conduct several training runs, each characterized by distinct settings for the policy's conservativeness (target KL $\in\{3,6,9,12\}$ for PPO, $\beta \in \{0.05,0.1,1,5\}$, $\alpha\in\{0.05,0.1,0.5,1\}$ for unlikelihood, and varying random seeds for preferred-FT), resulting in a total of 22 experimental runs. At every 100 training steps until models converge, policies are assessed on a suite of held-out prompts. For evaluation, we compute the mean reward according to the ground-truth reward metric and determine the average sequence-level KL\footnote{Defined as the sum of KL divergences at each time step.} relative to the reference policy, i.e., $\text{KL}\left(\pi\mid \mid \pi_\text{ref}\right)$. The empirical results indicate that DPO consistently establishes the most favorable frontier, achieving superior rewards at comparably low KL values. This finding is significant for several reasons. Primarily, although DPO and PPO both target the same optimization objective, DPO exhibits substantially greater efficiency, yielding a reward/KL trade-off that strictly outperforms PPO. Moreover, DPO's frontier surpasses that of PPO even in settings where PPO is granted access to true reward signals (PPO-GT).

\subsection{Can DPO scale to real preference datasets?}
\label{sec:dpo-real-datasets}
We proceed to assess the effectiveness of DPO fine-tuning when applied to tasks involving summarization and single-turn dialogue. For the summarization domain, established automatic metrics like ROUGE have been shown to correlate weakly with human judgments~\citep{stiennon2022learning}. Earlier studies suggest that language models fine-tuned with PPO based on human preference signals produce more preferred summaries. To evaluate various approaches, we generate model completions from the test partition of the TL;DR summarization dataset, then calculate the average win rate relative to reference completions within this set. Completions are sampled using different temperatures, ranging from 0.0 to 1.0; corresponding win rates are depicted in Figure~\ref{fig:frontier-tldr-main} (right). All methods, including DPO, PPO, and Preferred-FT, perform fine-tuning on an identical GPT-J SFT checkpoint\footnote{\url{https://huggingface.co/CarperAI/openai_summarize_tldr_sft}}. The results show that DPO reaches an approximate win rate of 61\% at temperature 0.0, outperforming PPO, which achieves around 57\% at its optimal sampling temperature (also 0.0). Additionally, DPO obtains a superior peak win rate when compared with the best-of-$N$ reference baseline. It should be emphasized that DPO's $\beta$ hyperparameter was not systematically optimized in these experiments, indicating these outcomes might not reflect the method’s full capability. Furthermore, our analysis demonstrates that DPO's performance remains relatively stable as the sampling temperature varies, in contrast to PPO, whose results can decline toward those of the base GPT-J model under higher temperatures. Preferred-FT exhibits negligible improvement over the SFT baseline. A direct comparison between DPO and PPO through human assessments, presented in Section~\ref{sec:human-judgments}, reveals that DPO samples drawn at temperature 0.25 were favored 58\% of the time over PPO samples collected at temperature 0.

For single-turn dialogue evaluation, we analyze various approaches using a segment of the test set from the Anthropic HH dataset \citep{bai2022training} containing single exchanges between humans and assistants. In the case of GPT-4, we assess each method by comparing its outputs to the human-preferred completions in the test set and calculating their respective win rates. Lacking a standard SFT baseline for this scenario, our process begins with a pre-trained Pythia-2.8B model. We apply Preferred-FT on selected completions to create a reference model that captures the completion distribution, which serves as the basis for subsequent DPO training. Our comparisons also include the best completion from 128 Preferred-FT samples—since the Best of $N$ metric saturates at 128 completions in our experiments (refer to Appendix Figure~\ref{fig:best-of-n})—as well as a 2-shot prompt configuration for the Pythia-2.8B model. In all cases, DPO matches or exceeds the performance of other methods at optimal sampling temperatures. 

Additionally, we evaluate an RLHF model, trained with PPO on the Anthropic HH dataset \footnote{\url{https://huggingface.co/reciprocate/ppo_hh_pythia-6B}} as distributed by a notable repository \footnote{\url{https://github.com/CarperAI/trlx/tree/main/examples/hh}}; however, no tested prompt or sampling temperature leads to better results than those achieved by the base Pythia-2.8B. Drawing from our findings with the TL;DR task, and recognizing that both PPO and Best of 128 target the same reward metric, we interpret Best of 128 as a useful benchmark for PPO-level performance. Overall, among the methods considered, DPO stands out for its computational efficiency and its consistent improvement over preferred completions in the Anthropic HH data, rivaling or surpassing the demanding Best of 128 baseline. As illustrated in Figure~\ref{fig:dialogue-main}, DPO also achieves its peak performance in relatively few training steps.

\subsection{Generalization to a new input distribution}

To more thoroughly assess how PPO and DPO handle shifts in input data, we test the policies obtained from our Reddit TL;DR summarization study on a novel data domain: news articles taken from the test split of the CNN/DailyMail dataset \citep{nallapati-etal-2016-abstractive}. For these evaluations, we utilize the top-performing sampling temperatures found in the original TL;DR experiment (0 and 0.25). The evaluation results are summarized in Table~\ref{tab:ood}. We calculate the GPT-4 win rate by comparing model-generated summaries to the reference summaries within each dataset, applying the same GPT-4 (C) prompt as for Reddit TL;DR—only modifying the wording from ``forum post'' to ``news article.'' On this new dataset, DPO maintains a considerable performance advantage over the PPO policy. These findings offer preliminary support that DPO-trained policies can generalize as effectively as those trained with PPO, even without the supplementary unlabeled Reddit TL;DR data that benefits PPO.

\subsection{Validating GPT-4 judgments with human judgments}
\label{sec:human-judgments}
To assess the accuracy of GPT-4’s evaluation in the context of the TL;DR summarization task, we conduct a human evaluation study using outputs from two separate GPT-4 prompts. The \textbf{GPT-4 (S)} (simple) prompt requests the rater to choose the summary that better captures the main information from the post. In contrast, the \textbf{GPT-4 (C)} (concise) prompt further asks which summary is more succinct. We focus on this additional dimension because, based on our observations, GPT-4 employing the \textbf{GPT-4 (S)} prompt tends to favor longer and more repetitive outputs than what is preferred by human raters. Detailed versions of both prompts are provided in Appendix~\ref{app:prompts}. We perform three comparative analyses: using the best-performing system (DPO, temp. 0.25), the lowest-performing (PPO, temp. 1.0), and an intermediate (SFT, temp. 0.25), in all cases benchmarking against PPO sampled greedily at its optimal temperature. The intent is to span a broad spectrum of output qualities; for DPO and PPO-1 comparisons, multiple human evaluations per comparison are collected (limited by available raters). Our results indicate that GPT-4's preferences align with human judgment nearly as closely as humans agree with each other, suggesting that GPT-4 constitutes a reliable proxy for human assessment. Notably, the \textbf{GPT-4 (C)} prompt yields win rates that more accurately mirror human preferences; hence, this prompt is adopted for the primary analysis in Section~\ref{sec:dpo-real-datasets}. Additional procedural details, including the rater interface and the roster of human participants, are available in Appendix~\ref{app:human-study}.

\section{Discussion}
Preference-based learning offers an effective and extensible approach for developing competent and aligned language models. In this work, we presented DPO, a streamlined training methodology that enables language models to learn from preference data without resorting to reinforcement learning techniques. Instead of transforming the preference learning problem into a conventional RL framework to apply standard RL algorithms, DPO constructs a relationship between language model policies and reward structures. This mapping allows for the direct training of language models to adhere to human preferences through a straightforward cross-entropy loss, eliminating the need for reinforcement learning or sacrificing model flexibility. Remarkably, DPO matches or surpasses the performance of prevalent RLHF approaches, including PPO-based methods, with minimal hyperparameter adjustment. This substantially lowers the barrier for incorporating human preferences in language model training.

\textbf{Limitations \& Future Work.} The findings prompt a number of key avenues for further investigation. How does the generalization capability of the DPO-trained policy on out-of-distribution samples compare to approaches that use explicitly defined reward functions? Preliminary observations indicate comparable generalization between DPO and PPO-based policies, yet a more thorough analysis is required. Another consideration is whether utilizing self-generated labels from the DPO policy can facilitate effective learning from unlabeled prompt data. Additionally, the mechanisms by which reward over-optimization appears in direct preference optimization—and whether the modest performance dip observed in Figure~\ref{fig:dialogue-main}-right is an example of this—warrant deeper exploration. Although current experiments involve models with up to 6B parameters, scaling DPO to state-of-the-art, significantly larger models presents a promising research trajectory. Concerning assessment practices, our results show that GPT-4-based win rates are sensitive to the prompting method, suggesting a need for research into prompt design to optimize automated system evaluations. Lastly, DPO’s methodology has potential applications that extend beyond language model preference learning, such as training generative models for other types of data. 

\end{document}